%----------------------------------------------------------------------
%%% APP: CONTENTS
%----------------------------------------------------------------------
% !TEX root = ./Main.tex


%The remainder of this paper is organized as follows. In \cref{app:organization}, we list a glossary of all the basic concepts and notation that we will utilize throughout the paper.
%In \cref{app:errata}, we expand the discussion for the related work in general-sum stochastic games and we describe all the typos and corrections in definitions or statements between the submission of main part and supplementary.
%In \cref{app:marl}, we prove all the properties of the Markov Decision Process (MDP), the value function and its gradient for the multi-agent reinforcement learning process when the model of episodic random stopping is in run. While similar machinery have been considered for the infinite horizon with $\gamma$ discounting regime and a standard analogy has been described in the literature, we prove here explicitly for the sake of completeness.
%In \cref{app:solutions}, we consider the different optimization solutions from variational perspective and their connection to Markov stationary Policy Nash Equilibria.
%In \cref{app:reinforce}, we conclude our preliminary results with the statistical guarantees of \reinforce policy gradient estimator.
%Finally, in \cref{app:asymptotic} we present the asymptotic convergence results for the variational stable points and\AG{I wrote we show  sos only for eager} establish the convergence rates for the ``eager'' and ``lazy'' stochastic policy gradient method for the second-order stationary points (See \cref{app:SOS}) and their deterministic counterparts (See \cref{app:deterministic}).

\begin{table}[tbp]
\label{tab:index}
\centering
\renewcommand{\arraystretch}{1.2}
\begin{tabular}{cp{12cm}}
\toprule
\textsc{Notation}
	&\textsc{Description}
	\\ % <-- added & and content for each column
\midrule
$\rstate\in\states$
	&States of the game
	\\
$\pure_\play\in\pures_\play$
	& Actions of agent $\play\in\players$
	\\
%		\hline
$\traj$
	&Episode trajectory
	\\
%		\hline
$\stopTime(\traj)$
	&Episode stopping time
	\\
%		\hline
$\stopPr$
	&Minimum stopping probability
	\\
%		\hline
$\realReward_{\play,\time}$
	&Realized reward of $\play$-th player at time $\time$
	\\
%		\hline
$\curr[\step]$
	&Step size at episode $\run$
	\\
%		\hline
$\expar_\run$
	&Explicit exploration parameter at episode $\run$
	\\
%		\hline
$\gvalue(\curr)$
	&Policy gradients at policy $\curr$ of episode $\run$
%	When we write $\gvalue_{\play,\pure_\play,\rstate}(\policy)$, we refer to the gradient value of agent $\play\in\players$ for the specific pure action $\pure_\play\in\pures_\play$ and state $\rstate\in\states$ evaluated in the policy $\policy$.
	\\
%		\hline
$\signal_\run$
	&Policy gradient proxy at episode $\run$.
%	We use the same notation with $\gvalue_\run(\policy)$
	\\
%		\hline
%$\dstate_\run$
%	&Aggregate gradients at round $\run$
%	\\
\bottomrule
\end{tabular}
\caption{Refresher of the most common notations used in our paper.}
\end{table}